\hypertarget{introduction}{%
\section{1. Introduction}\label{introduction}}

SOAP~\citep{vyas2024soap} applies Adam-style
adaptation in the coordinate system of matrix preconditioners. Computing
and maintaining that coordinate system adds work beyond diagonal
optimizers. A practical approximation must address both its numerical
history and how that history is stored and executed on a GPU.

We investigate one concrete implementation point: full dense factors and
two-sided matrix coverage, BF16 persistent state, a fixed mixture of
cheaper dense tracking steps and periodic QR refreshes, and compilation
of the complete matrix-update regions. The objective is to document how
closely this fixed implementation follows an audited SOAP reference on
an affordable language modeling workload, together with the measured
time and memory tradeoff.

The evidence has three parts. First, two initialization seeds compare
the declared 10,000-call endpoint against stock-beta QR10. Second, short
full-BF16 QR20 and QR50 controls test whether intermediate tracking
provides useful learning behavior relative to simpler refresh schedules.
Third, a fixed-input profile separates redundant products in QR
refinement from the solver cost. The policy, precision rules, production
optimizer hashes, and shared training-worker hash remain fixed across
the two main candidate runs; campaign-specific runners are separately
pinned.

Our contribution is the specified dense tracking/repair/precision
implementation and its reproducible case study. The learning curves
retain a measurable loss gap, especially early in training. The timing
comparison includes storage, batching, arithmetic, compilation, and
cadence changes relative to the original reference. These components
should remain visible when interpreting efficiency.
